
\documentclass{article}
\usepackage{colortbl}
\usepackage{makecell}
\usepackage{multirow}
\usepackage{supertabular}

\begin{document}

\newcounter{utterance}

\centering \large Interaction Transcript for game `eqbench', experiment `eq\_bench', episode 112 with reflexion{-}Qwen3.5{-}27B.
\vspace{24pt}

{ \footnotesize  \setcounter{utterance}{1}
\setlength{\tabcolsep}{0pt}
\begin{supertabular}{c@{$\;$}|p{.15\linewidth}@{}p{.15\linewidth}p{.15\linewidth}p{.15\linewidth}p{.15\linewidth}p{.15\linewidth}}
    \# & \multicolumn{2}{c}{Player} && \multicolumn{2}{c}{Game Master} \\
    \hline

    \theutterance \stepcounter{utterance}  
    & & & \multicolumn{4}{p{0.6\linewidth}}{
        \cellcolor[rgb]{0.9,0.9,0.9}{
            \makecell[{{p{\linewidth}}}]{
                \texttt{\tiny{[P1$\langle$GM]}}
                \texttt{You have attempted this same task before and did not succeed. These are the plans you wrote after those attempts.} \\
\texttt{Trial \#1: New plan: My previous attempt failed because I assigned a non{-}zero score to "Angry" (1) and "Dismissed" (2), which likely contradicted the specific emotional nuance of the dialogue where Ben is playfully engaging rather than genuinely offended or ignored. Ben's response ("Guilty as charged, sweetheart") indicates he accepts the challenge and enjoys the banter, making "Challenged" the dominant emotion while the others should be strictly zero. Next time, I will analyze the tone of the dialogue more critically to ensure that only the most logically supported emotion receives a high score, and all other plausible but incorrect emotions are assigned a score of 0 to avoid diluting the primary emotional state. Specifically, I will set "Challenged" to a high value (e.g., 9 or 10) and explicitly set "Dismissed", "Angry", and "Apathetic" to 0.} \\
\\ 
\texttt{Trial \#2: New plan: My previous attempt failed because I assigned a score of 0 to "Dismissed," overlooking the nuance in Ben's opening line ("Well, well, well, if it isn't the Ice Queen herself") which is a direct response to being previously ignored or coldly received. While "Challenged" is the dominant emotion due to the competitive banter, "Dismissed" is a valid secondary emotion reflecting the context of the "Ice Queen" persona. I incorrectly assumed a single{-}emotion model where only one score could be non{-}zero. Next time, I will analyze the dialogue for layered emotional responses. I will assign a high score (e.g., 9) to "Challenged" as the primary reaction to the "game of skill" comment, but I will also assign a moderate score (e.g., 4) to "Dismissed" to account for the initial cold reception implied by the "Ice Queen" label, while keeping "Angry" and "Apathetic" at 0. This ensures the scoring reflects the complexity of the interaction rather than forcing a single{-}dimension interpretation.} \\
\\ 
\\ 
\\ 
\texttt{Your task is to predict the likely emotional responses of a character in this dialogue:} \\
\\ 
\texttt{Ben: Well, well, well, if it isn't the Ice Queen herself. I've heard tales of your cold{-}blooded poker face.} \\
\texttt{Esther: And you must be Benny the Bluff. I've heard you're quite the charmer. } \\
\texttt{Ben: Guilty as charged, sweetheart. But don't let that fool you. I've got a few tricks up my sleeve.} \\
\texttt{Esther: I'm sure you do. But remember, Ben, this isn't a magic show. It's a game of skill. } \\
\texttt{{[}End dialogue{]}} \\
\\ 
\texttt{At the end of this dialogue, Ben would feel...} \\
\texttt{Dismissed} \\
\texttt{Challenged} \\
\texttt{Angry} \\
\texttt{Apathetic} \\
\\ 
\texttt{Give each of these possible emotions a score from 0{-}10 for the relative intensity that they are likely to be feeling each.} \\
\\ 
\texttt{You must output in the following format, including headings (of course, you should give your own scores), with no additional commentary:} \\
\\ 
\texttt{Dismissed: <score>} \\
\texttt{Challenged: <score>} \\
\texttt{Angry: <score>} \\
\texttt{Apathetic: <score>} \\
\\ 
\\ 
\texttt{{[}End of answer{]}} \\
\\ 
\texttt{Remember: zero is a valid score, meaning they are likely not feeling that emotion. You must score at least one emotion > 0.} \\
\\ 
\texttt{Your answer:} \\
            }
        }
    }
     \\ \\

    \theutterance \stepcounter{utterance}  
    & \multicolumn{4}{p{0.6\linewidth}}{
        \cellcolor[rgb]{0.561,0.737,0.561}{
            \makecell[{{p{\linewidth}}}]{
                \texttt{\tiny{[P1$\rangle$GM]}}
                \texttt{Dismissed: 0} \\
\texttt{Challenged: 9} \\
\texttt{Angry: 0} \\
\texttt{Apathetic: 0} \\
            }
        }
    }
    & & \\ \\

    \theutterance \stepcounter{utterance}  
    & & & \multicolumn{2}{p{0.3\linewidth}}{
        \cellcolor[rgb]{0.9,0.9,0.9}{
            \makecell[{{p{\linewidth}}}]{
                \texttt{\tiny{[GM$|$GM]}}
                \texttt{Dismissed: 0} \\
\texttt{Challenged: 9} \\
\texttt{Angry: 0} \\
\texttt{Apathetic: 0} \\
            }
        }
    }
    & & \\ \\

    \theutterance \stepcounter{utterance}  
    & & & \multicolumn{2}{p{0.3\linewidth}}{
        \cellcolor[rgb]{0.9,0.9,0.9}{
            \makecell[{{p{\linewidth}}}]{
                \texttt{\tiny{[GM$|$GM]}}
                \texttt{Dismissed: 3} \\
\texttt{Challenged: 8} \\
\texttt{Angry: 0} \\
\texttt{Apathetic: 0} \\
            }
        }
    }
    & & \\ \\

    \theutterance \stepcounter{utterance}  
    & & & \multicolumn{2}{p{0.3\linewidth}}{
        \cellcolor[rgb]{0.9,0.9,0.9}{
            \makecell[{{p{\linewidth}}}]{
                \texttt{\tiny{[GM$|$GM]}}
                \texttt{game\_result = LOSE} \\
            }
        }
    }
    & & \\ \\

    \theutterance \stepcounter{utterance}  
    & & \multicolumn{4}{p{0.6\linewidth}}{
        \cellcolor[rgb]{0.95,0.95,0.95}{
            \makecell[{{p{\linewidth}}}]{
                \texttt{\tiny{[P1$\leftrightarrow$P]}}
                \texttt{Post{-}episode Reflexion (not sent to GM):} \\
\texttt{New plan: My previous attempt failed because I assigned a score of 0 to "Dismissed," overlooking the nuance in Ben's opening line ("Well, well, well, if it isn't the Ice Queen herself") which is a direct response to being previously ignored or coldly received. While "Challenged" is the dominant emotion due to the competitive banter, "Dismissed" is a valid secondary emotion reflecting the context of the "Ice Queen" persona. I incorrectly assumed a single{-}emotion model where only one score could be non{-}zero. Next time, I will analyze the dialogue for layered emotional responses. I will assign a high score (e.g., 9) to "Challenged" as the primary reaction to the "game of skill" comment, but I will also assign a moderate score (e.g., 4) to "Dismissed" to account for the initial cold reception implied by the "Ice Queen" label, while keeping "Angry" and "Apathetic" at 0. This ensures the scoring reflects the complexity of the interaction rather than forcing a single{-}dimension interpretation.} \\
            }
        }
    }
    & \\ \\

\end{supertabular}
}

\end{document}
